\section{Two finite systems of idempotents}
\label{sec:idempotents}

The finite diagrams used in our construction are indexed by idempotent
operators over $\F_2$ and $\F_3$.  Their essential property is elementary:
an additive assignment to either diagram forces the value at the identity
to be annihilated by the corresponding prime.  These diagrams and their
additivity argument are shared with the
companion~\cite[Section~2.1, Lemma~2.1]{OpenAIOvergroups}.
We give the complete argument before realizing them by homomorphisms of
groups.

For $\ell\in\{2,3\}$, let $D_\ell$ be the unital $\F_\ell$-algebra of
operators on the free vector space with basis $\{0,1\}^{\N}$ generated by
$l_0,l_1,r_0,r_1$.  Here $l_i$ prefixes the symbol $i$, and $r_i$ deletes
an initial $i$, sending basis vectors with the other initial symbol to
zero.  Writing $I$ for the identity operator, we have
\begin{equation}\label{eq:binary-prefix-relations}
 r_i l_j=\delta_{ij}I,
 \qquad l_0r_0+l_1r_1=I.
\end{equation}
For a finite binary word $v$, let $l_v$ prefix $v$, let $r_v$ delete $v$
when it occurs as a prefix and give zero otherwise, and put $E_v=l_vr_v$.
Thus $E_v$ projects onto the span of the sequences beginning with $v$.

If $(v_1,v_2,v_3)$ is a prefix-free list, then
\[
 (l_{v_i}r_{v_j})(l_{v_k}r_{v_m})
   =\delta_{jk}l_{v_i}r_{v_m}.
\]
These operators are nonzero matrix units, so their span is a copy of
$M_3(\F_\ell)$.  We use three such algebras, connected by shared
diagonal projections.  Write
\[
 A=E_{00},\qquad A_0=E_{000},\qquad A_1=E_{001},
 \qquad B'=E_{01},\qquad C'=E_1.
\]
Their defining lists, diagonal projections, and identity operators are
\begin{equation}\label{eq:three-prefix-lists}
\begin{array}{c|c|c}
 \text{prefix list} & \text{diagonal projections} & \text{identity}\\ \hline
 (00,01,1)  & A,B',C'   & I\\
 (000,01,1) & A_0,B',C' & A_0+B'+C'\\
 (001,01,1) & A_1,B',C' & A_1+B'+C'.
\end{array}
\end{equation}
The last two identities are not $I$, while
$A=A_0+A_1$ and $A_0A_1=A_1A_0=0$.  Let $\mathcal P_\ell$ be the union,
as a set of actual
operators in $D_\ell$, of all idempotents in these three matrix
algebras.  It is finite and contains $I$.  An operator belonging to two
of the algebras has just one index in $\mathcal P_\ell$.

Two operators $E,P$ are called \emph{orthogonal} if $EP=PE=0$.
The relations below include all orthogonal pairs whose sum belongs to
$\mathcal P_\ell$, even when the pair does not belong to a single one of
the three matrix algebras.  In particular, they include
$A=A_0+A_1$, whose summands come from the two smaller support algebras.

\begin{lemma}\label{lem:idempotents}
Let $\ell\in\{2,3\}$, let $T$ be an abelian group, and let
$(z_E)_{E\in\mathcal P_\ell}$ be elements of $T$.  Suppose that
\begin{equation}\label{eq:idempotent-additivity}
 z_{E+P}=z_E+z_P
\end{equation}
whenever $E,P,E+P\in\mathcal P_\ell$ and $EP=PE=0$.
Then $\ell z_I=0$.
\end{lemma}

\begin{proof}
We first compare the values of rank-one idempotents within any one of
the three matrix algebras.  For a column vector $u\in\F_\ell^3$ and a row
covector $\lambda$ with $\lambda u=1$, write
\[
 x(u,\lambda)=z_{u\lambda}.
\]
Fix a nonzero vector $v\in\ker\lambda$.  We claim that the difference
\begin{equation}\label{eq:fixed-increment}
 x(u+v,\lambda)-x(u,\lambda)
\end{equation}
does not depend on the vector $u$ subject to $\lambda u=1$.

Take another such vector $u'$.  Suppose first that $u'-u$ and $v$ are
linearly independent.  They form a basis of $\ker\lambda$, and hence
$u,u'-u,v$ form a basis of $\F_\ell^3$.  There is consequently a covector
$\mu$ satisfying
\[
 \mu u=\mu u'=0,\qquad \mu v=1.
\]
The identity
\[
 u\lambda+v\mu=(u+v)\lambda+v(\mu-\lambda)
\]
expresses one rank-two idempotent as two orthogonal sums of rank-one
idempotents.  Indeed, $\lambda v=\mu u=0$ on the left, while
$\lambda v=(\mu-\lambda)(u+v)=0$ on the right; all four rank-one
factors have their defining covector equal to one on their defining
vector.  Applying \eqref{eq:idempotent-additivity} gives
\begin{equation}\label{eq:rank-one-difference}
 x(u+v,\lambda)-x(u,\lambda)
   =x(v,\mu)-x(v,\mu-\lambda).
\end{equation}
The same covector $\mu$ gives the same right-hand side for $u'$, proving
the claim in this case.

If $u'-u$ lies in the span of $v$, choose
$w\in\ker\lambda\setminus\langle v\rangle$ and put $u''=u+w$.
Both $u''-u$ and $u''-u'$ are independent of $v$.  Comparing first $u$
with $u''$ and then $u'$ with $u''$ proves the claim in general.
The kernel of $\lambda$ has dimension two, so this argument also works
over $\F_2$; no choice of a third nonzero scalar is being used.

Apply the claim to the successive vectors
$u,u+v,\ldots,u+(\ell-1)v$.  The $\ell$ successive differences are
equal, and they sum to zero because $\ell v=0$.  Therefore
\begin{equation}\label{eq:rank-one-prime-invariance}
 \ell x(u+v,\lambda)=\ell x(u,\lambda).
\end{equation}
The same conclusion is immediate for $v=0$, so $\ell x(u,\lambda)$
is independent of $u$ whenever $\lambda$ is fixed and $\lambda u=1$.
For the covector comparison, work within this same matrix algebra
and define a new assignment $z'_E=z_{E^{\mathsf T}}$.
Transposition preserves idempotence and orthogonal sums, so $z'$ obeys
the same additivity relations.  Apply the invariance just proved to
$x'(v,\mu)=z'_{v\mu}=z_{\mu^{\mathsf T}v^{\mathsf T}}$.
Fixing $\mu$ and varying $v$ now shows that $\ell x(u,\lambda)$
is independent of $\lambda$ whenever $u$ is fixed and $\lambda u=1$.
This does not assume that $z_E=z_{E^{\mathsf T}}$.

In particular the three diagonal matrix units have equal
$\ell$-multiples.  To see the comparison explicitly, for distinct basis
vectors $e_i,e_j$ take $\lambda=e_i^*+e_j^*$.  First change the covector
at $e_i$ from $e_i^*$ to $\lambda$, then change the vector from $e_i$
to $e_j$, and finally change the covector to $e_j^*$.  Each change
preserves the $\ell$-multiple, giving
\[
 \ell z_{e_i e_i^*}=\ell z_{e_j e_j^*}.
\]

It remains to use the shared projections displayed above.  The
comparison of diagonal units in each algebra gives
\[
 \ell z_A=\ell z_{B'}=\ell z_{C'}
           =\ell z_{A_0}=\ell z_{A_1}.
\]
Call this common value $q\in T$.  The projection identity
$A=A_0+A_1$ is an orthogonal sum across the two smaller support
algebras.  Its three operators belong to $\mathcal P_\ell$, so
\eqref{eq:idempotent-additivity} gives $q=q+q$.  Thus $q=0$.
Finally, $A+B'$ belongs to the first matrix algebra, and
$I=(A+B')+C'$.  Applying additivity twice yields
\[
 \ell z_I=\ell z_A+\ell z_{B'}+\ell z_{C'}=0,
\]
as required.
\end{proof}
